% ============================================================

\begin{frame}{Same National Shift, Unequal Local Exposure}
\hypertarget{src:variation}{}
\small
Estimating the effect of the adversarial arena on a race requires variation in litigation that local politics does not manufacture: municipalities that litigate more are not a random sample of the country.
\smallskip

National, subject-level litigation trends land unevenly across municipalities, according to each one's \emph{prior} caseload mix. Two municipalities meet the same national shift with different exposure:
\vspace{2pt}
\begin{center}
\begin{tikzpicture}[
  font=\footnotesize,
  wave/.style={draw=myblue, thick, rounded corners, fill=myblue!8, align=center, inner sep=5pt},
  muni/.style={draw=black!55, rounded corners, align=center, inner sep=4pt, minimum width=3.3cm},
  hi/.style={-{Latex}, myblue, line width=2.6pt},
  lo/.style={-{Latex}, myblue!45, line width=0.9pt},
]
  \node[wave] (nat) at (0,2.25) {\textbf{National shift} in litigation topic $k$\\ \scriptsize (e.g.\ \emph{fraude \`a cota}), measured leave-own-state-out};
  \node[muni, fill=myblue!16] (a) at (-3.5,0) {\textbf{Muni A}\\ \scriptsize \emph{heavy} in $k$ in 2020};
  \node[muni, fill=myblue!4]  (b) at (3.5,0)  {\textbf{Muni B}\\ \scriptsize \emph{light} in $k$ in 2020};
  \draw[hi] (nat) -- (a) node[midway,above left=-2pt]{large push};
  \draw[lo] (nat) -- (b) node[midway,above right=-2pt]{small push};
  \node[below=3pt of a, myblue]   {\scriptsize \textbf{high} predicted exposure $Z_A$};
  \node[below=3pt of b, black!55] {\scriptsize \textbf{low} predicted exposure $Z_B$};
\end{tikzpicture}
\end{center}
\vspace{2pt}
\centerline{\footnotesize The comparison is between municipalities pushed harder by \emph{identical} national trends.}
\smallskip
\footnotesize Two ingredients construct it: each municipality's 2020 caseload mix across litigation subjects, and the national shift in each subject measured leave-own-state-out. Both rest on a prior decision about which filings count as adversarial.
\end{frame}
